\section{Three entries in the original period and height estimate}

% BEGIN P88A1
\paragraph{P88A1. One constant and one threshold for the completed entries.}
Write $\mathcal C=\mathrm{WithChN}(\mathcal A)$ for the class
in X87S4. For the conjunction of entries three, five and eight of
$\mathrm{LinkDIV}_{\rm at}$ and the rationality clause of entry
four, the source's outer closed-family statement holds on
$\mathrm{Des}(\mathcal C,3)$. Take $\epsilon_1=1/53$.
For each $0<\epsilon\le\epsilon_1$ choose the zero transport
parameter before the family or its height constant. Let
$\mathcal F$ belong to the class and satisfy the prescribed
largeness condition. Keep the level-prime and stabilizer-size
bounds in the original conditional statement. Both are also
available by projection from the design membership. The following
conclusions use the very datum $D=\mathcal F.D(x,i)$ and size
record $L=\mathcal F.L(x,i)$.

By H87P1 the support assertions hold for $x\ge\mathcal F.x_1+1$.
On precisely that guarded range, X87S1 identifies the selected
bundle and its norm with the linked point. R88L2 gives the full
row lattice on the original sharp side and also on the base side.
R88L3 supplies the rational period pairing at the recorded
coefficient embedding. H88E5 gives a constant $C$ and an
eventual threshold for which, for every word $u$ and allowed
lift $j$,
\[
 \left|\sum_{w\in L.\mathrm{supp}}2\ell_w e_w E_{u,j,w}
       -\left(\sum_{s\in TT}a_s\right)\log(g_0g_1)\right|
    \le C B_0 .
\]
Here $\ell_w=\log(\mathrm{N}w)/[K:\mathbf Q]$ for the coefficient
field $K$, $e_w$ is the local ramification index and
$B_0=fmbx^\epsilon$. H87P2 uses these very weights on the two
distinguished prime fibres. Use the maximum of the preceding
thresholds. This proves all four assertions for one family $C,x_0$,
followed by all $x\ge x_0$, all indices $i$, and all the row and
lift choices in those assertions. The same conclusion follows
under the additional $\mathrm{LinkEta}$ and $\mathrm{DIVinN}$
hypotheses of the period and height contract.

The full contract requires the common unitary comparison in
entry four, the integral map and frame bound in entries one and
two, the two word-height bounds in entries six and seven, and
the comparison between the corrected heights of two words in
entry nine. In assembling all nine entries, bounds with a
variable constant may be combined by taking a maximum when
their multiplying allowances are nonnegative; the exact
identities, fixed frame bound and fixed map-height bound retain
their own stated scopes.
% END P88A1
